Digitization of hand drawn circuit schematics has remained challenging because reliable interpretation depends on both symbol localization and recovery of thin connecting wires. Recent progress has been reported in handwritten and engineering diagram recognition through detector based pipelines, graph enhanced postprocessing, and domain specific datasets \cite{Thoma2021,Rachala2022,Shen2023,Schaefer2021,Elyan2020,Lee2020,Yang2024,Hemker2024,MorenoGarcia2024}. In most of these studies, strong performance has been obtained for component or symbol detection, whereas wire or pipeline recovery has remained sensitive to local clutter, rotation, scan noise, and low contrast traces. The resulting gap has limited the direct conversion of offline sketches into simulation ready circuit structure.

The document analysis literature has shown that binarization and line recovery remain decisive when fine strokes must be preserved. Classical thresholding methods such as Otsu, Niblack, and Sauvola remain widely used because they are interpretable and easy to control \cite{Otsu1979,Niblack1986,Sauvola2000}. Their limitations on degraded documents have been studied extensively, and improved adaptive strategies have been proposed through contrast based thresholding, local background estimation, and efficient integral image implementations \cite{Gatos2006,Shafait2008,Su2013}. At the same time, skeletonization and line extraction have been developed into mature tools for preserving topology and tracing elongated structures \cite{ZhangSuen1984,Lam1992,DudaHart1972,Gioi2010}. These ideas are especially relevant when the objective is not merely foreground segmentation, but recovery of sparse geometric structure from weak and fragmented strokes.

Related work in engineering drawings has also indicated that symbol recognition alone is insufficient. Deep detectors have been used successfully for symbols in piping and instrumentation diagrams, single line electrical drawings, and related symbolic sketches, yet downstream structural inference still requires domain aware processing of lines, ports, and connectivity \cite{Elyan2020,Lee2020,Yang2024,Hemker2024,MorenoGarcia2024}. Hand drawn circuit studies have similarly relied on object detection followed by node or line reasoning, often through Hough based or heuristic geometric procedures \cite{Rachala2022}. These observations suggest that hybrid methods remain attractive when wire specific labeled data are limited.

In the earlier manuscript version, a hybrid strategy was adopted in which component detections were used as priors while wire extraction was performed by a classical pipeline based on Sauvola thresholding, connected component analysis, and segment merging. That formulation established the value of component guided preprocessing, but the strongest validated result from the current repository differs materially from the earlier draft. The merge based wire result of 0.647 no longer reflects the best reproducible method. A reproduced reference pipeline now yields an F1 score of 0.7066, and the strongest validated method in the present work yields an F1 score of 0.8215 on the same 23 image benchmark.

The present revision has therefore been focused on the wire method itself. The key observation was that line fitting directly from connected components compresses image structure too early. Fragmented traces, redundant segments, and ambiguous endpoints were repeatedly produced in dense regions. To address this issue, a skeleton graph representation was introduced. Adaptive threshold masks were converted into graph paths, candidate wire hypotheses were scored by structural and component guided cues, and a topology aware selection step was applied to preserve a sparse and plausible explanation of the circuit layout.

The resulting manuscript has been written as a MethodsX article because the principal contribution is methodological. A reproducible hybrid wire extraction procedure is documented, its behavior is validated on a manually annotated benchmark, and its design choices are presented in a form that can be inspected and reused in related document analysis tasks.
